\section{Evaluation} \label{sec:evaluation}
This section presents the evaluation of the Rocq-based proof checker and compares it with the Imandra proof checker described in \cite{certified_proof_checker}. The evaluation is structured along two complementary dimensions: capability and maintainability.

\subsection{Capability Evaluation}

\autoref{tab:capability-comparison} compares the functional capabilities of the Rocq proof checker against the Imandra baseline. Each feature is marked as supported (\checkmark) or not supported (---).

\begin{table}[h]
\centering
\begin{tabularx}{\textwidth}{l c c}
\toprule
\textbf{Feature} & \textbf{Imandra} & \textbf{Rocq} \\
\midrule
ReLU activation support & \checkmark & \checkmark \\
Identity activation support & \checkmark & \checkmark \\
Theory lemmas support & \checkmark & --- \\
Theory lemmas verification & --- & --- \\
Proof Checker & \checkmark & \checkmark \\
Proof Checker Soundness & \checkmark & \checkmark \\
Executable Proof Checker & \checkmark & --- \\
\bottomrule
\end{tabularx}
\caption{Capability comparison between the Imandra and Rocq proof checkers.}
\label{tab:capability-comparison}
\end{table}

TODO: DISCUSS THE CAPABILITIES TABLE
% The capability comparison shows that the Rocq proof checker achieves functional parity with the Imandra baseline on the core verification features: both support ReLU activation constraints, both include a complete proof checker, and both are accompanied by a formal soundness theorem. The only unimplemented feature is the verification of theory lemmas, which in the Imandra checker are used to validate intermediate bound-tightening deductions within the proof tree. The Rocq checker currently propagates bounds through case splits but does not independently verify the lemma steps that Marabou may include at internal nodes.


\subsection{Maintainability Evaluation}
The maintainability evaluation analyses the codebases of both implementations to assess their relative maintainability. The following metrics will be used to ensure a fair comparison:

\begin{itemize}
    \item \textbf{Lines of code.} The total number of lines of code in each implementation, broken down by file, to assess the overall size and complexity of the codebase.
    \item \textbf{Number of auxiliary lemmas.} The number of helper lemmas required to support the main proof is used as a proxy for proof complexity.
    \item \textbf{External library dependencies.} The number and nature of external libraries relied upon by each implementation.
\end{itemize}

TODO WRITE EVALUATION RESULTS.
